% Seminar note temvlate
% Usage: vlace your content in the sections below and run `vdflatex seminar_temvlate.tex`
\documentclass[11vt,a4vaver]{article}
\usepackage[utf8]{inputenc}
\usepackage[T1]{fontenc}
\usepackage{lmodern}
\usepackage[margin=2.5cm]{geometry}
\usepackage{microtype}
\usepackage{amsmath,amssymb,amsthm}
\usepackage{graphicx}
\usepackage{caption}
\usepackage{enumitem}
\usepackage{hyperref}
\usepackage{bookmark}
\usepackage{fancyhdr}

% Header / footer
\pagestyle{fancy}
\fancyhf{}
\fancyhead[L]{Seminar notes}
\fancyhead[R]{\leftmark}
\fancyfoot[C]{\thepage}

% Theorem-like environments (optional)
\newtheorem{theorem}{Theorem}[section]
\newtheorem{lemma}[theorem]{Lemma}
\theoremstyle{definition}
\newtheorem{definition}[theorem]{Definition}

% Document info - edit these
\title{Some Discussion for a tyve of Generalized Heisenberg Lie Algebra}
\author{Hanfeng Huang}
\date{\today}

\begin{document}

\maketitle

\begin{abstract}
% 摘要：簡短描述本次報告/讀書會的核心內容與結論（3–6 行）。
In this note, we study a class of step two nilpotent Lie algebras arising
from a linear space equipped with a constant Poisson structure, which could 
be realized as a generalized Heisenberg Lie algebra. We discuss the structure of these algebras, 
associating the singular point of its coadjoint orbit to the affine variety.

\end{abstract}

% 目錄
\tableofcontents
\vspace{1em}

% 建議結構：Introduction -> Main Sections -> Discussion -> References
\section{Preliminary}
\label{sec:intro}
% 背景、問題動機、講者與議題概述
In this section, we briefly review some basic notion of Poisson manifolds, and discuss some 
properties of constant Poisson structure on a finite dimensional linear space.

\subsection{Poisson Manifold and Hamiltonian vector field}

\begin{definition}[Poisson       manifold]
    A Poisson manifold $(\mathcal{M}, \{\cdot,\cdot\})$is a smooth manifold $\mathcal{M} $ equipped with a 
    bilinear operation $\{\cdot,\cdot\}: \mathcal{C} ^\infty(\mathcal{M} ) \times \mathcal{C} ^\infty(\mathcal{M} ) 
    \to \mathcal{C} ^\infty(\mathcal{M} )$, called the Poisson bracket, satisfying the following properties:
    \begin{enumerate}
        \item Skew-symmetry: $\{f,g\} = -\{g,f\}$ for all $f,g \in \mathcal{C} ^\infty(\mathcal{M} )$.
        \item Jacobi identity: $\{f,\{g,h\}\} + \{g,\{h,f\}\} + \{h,\{f,g\}\} = 0$ for all $f,g,h \in \mathcal{C} ^\infty(\mathcal{M} )$.
        \item Leibniz rule: $\{f,gh\} = g\{f,h\} + h\{f,g\}$ for all $f,g,h \in \mathcal{C} ^\infty(\mathcal{M} )$.
    \end{enumerate}
\end{definition}



\begin{definition}[Hamiltonian vector field]
    Given a Poisson manifold $(\mathcal{M}, \{\cdot,\cdot\})$ and a smooth function $H \in \mathcal{C} ^\infty(\mathcal{M} )$, 
    the Hamiltonian vector field $X_H$ associated with $H$ is defined by the relation
    \[
        <df, X_H> = \{f,H\}, \quad \forall f \in \mathcal{C} ^\infty(\mathcal{M} ),
    \]
    where $<\cdot,\cdot>$ denotes the natural pairing between differential forms and vector fields.
\end{definition}

\subsection{Constant Poisson structure on a linear space}

\begin{definition}[Constant Poisson structure]
    Let $V$ be a finite dimensional linear space over $\mathbb{R}$ , 
    and let $\{\cdot,\cdot\}$ be a Poisson bracket on $V$. 
    We say that the Poisson structure is constant if the 
    Poisson bracket of any two linear functions on $V$ is constant. 
    In other words, for any linear functions 
    $\xi,\eta\in V^*$, we have $\{\xi,\eta\} \in \mathbb{R}$.
\end{definition}

Since the bracket with any constant function is zero, 
the bracket actually defines a skew-symmetric bilinear form $\Omega$
on the dual space $V^*$ by $\Omega(\xi,\eta) := \{\eta,\xi\}$. 
One can check that the Jacobi identity is automatically satisfied for a constant Poisson structure,
and thus the Poisson         bracket is completely determined 
by the skew-symmetric bilinear form $\Omega$ on $V^*$ once we identify the tangent space $T_vV$
of $V$ at any point $v\in V$ with $V$ itself and its double dual $V^{**}$, and also
the differential $d_v\xi$ of any linear function $\xi\in V^*$ with $\xi$
itself. 
Note that 
the Hamiltonian vector field $X_\xi$  is given by
$\Omega_\sharp(\xi):=\iota_{X_\xi}\Omega\in V$, 
which is a constant vector field on $V$.

Next we restrict our attention to the space of all symmetric 
bilinear forms on $V$, denoted by $\text{S}^2(V)$. For $\sigma \in
\text{S}^2(V)$, write $\sigma(v)=\frac{1}{2}\sigma(v,v)$ as a smooth 
function on $V$. Note that $d_v\sigma= \iota_v \sigma\in T^*_vV 
\cong V^*$. 

We first compute the bracket of two quadratic $\sigma_1,\sigma_2$:
\begin{align*}
    \{\sigma_1,\sigma_2\}(v)&=<d_v{\sigma_1},X_{\sigma_2}(v)>=
    <\iota_v\sigma_1,X_{\sigma_2}(v)>=\{\iota_v\sigma_1,\sigma_2\}(v)\\
    &=-\{\sigma_2,\iota_v\sigma_1\}(v)=-<d_v\sigma_2,X_{\iota_v\sigma_1}(v)>=
    -<\iota_v\sigma_2,X_{\iota_v\sigma_1}(v)>\\
    &=-\{\iota_v\sigma_2,\iota_v\sigma_1\}=\{\iota_v\sigma_1,\iota_v\sigma_2\}
        \end{align*}
    

We first compute the bracket of a quadratic function $\sigma$ and a 
linear function $\xi$:
\begin{align}
    \{\sigma,\xi\}(v) &= <d_v\sigma, X_\xi>=\sigma(v,X_\xi)=<\iota_{X_\xi}\sigma, v> .
    \end{align} 
    The formula above shows that the bracket $\{\sigma,\xi\}$ 
    is a linear function.






\end{document}
